%% =====================================================================
%%  T-28-02  The Goal Line
%%  -------------------------------------------------------------------
%%  One idea per file. Core is mandatory. Slots in canonical order.
%%  Max 700 words. No layout commands. No filler. New source material
%%  for this entry goes in sources/B3/T-28-02.tex -- never by
%%  rewriting this file (docs/RULES.md R0.1).
%%  =====================================================================
%% @kind: topic
%% @id: T-28-02
%% @title: The Goal Line
%% @subtitle: Set the target, reach it, stop there
%% @chapter: c28
%% @order: 20
%% @status: stable
%% @sources: B3
%% @updated: 2026-09-19

\topic{T-28-02}{The Goal Line}{Set the target, reach it, stop there}

\begin{core}
The most dangerous moment of any campaign is the one in which it succeeds. Victory intoxicates; intoxication re-aims the instrument at targets nobody chose; and pressing past the chosen mark manufactures more enemies than the win defeated. Self-mastery's final test is the discipline to stop where you decided to stop.
\end{core}
\begin{mechanism}
B3 states it as a rule of timing: the hour of triumph is frequently the hour of maximum danger, because arrogance and overconfidence push you past the goal you had aimed at, and by going too far you manufacture enemies out of spectators. Success is uniquely persuasive precisely because it was earned --- the operator who just beat the odds believes in their momentum, and momentum has no eyes. The law's story is Cyrus: king of the world after Babylon, warned by the queen he marched against --- be content, keep your forces whole, rule what you won --- and he crossed one river too many, to end, by her oath, drinking the blood he had demanded. No cleverness substitutes for the line itself: fix the goal in advance, and on reaching it, halt. The line works because it was set by the planner rather than the reveller; a goal chosen mid-victory is already the intoxication speaking.
\end{mechanism}
\begin{conditions}
Strongest at actual victories --- the test follows success, never failure --- and visible in miniature whenever a negotiator keeps selling after the yes. It presumes the target was set in advance; without a pre-set mark there is nothing to stop at, only drift. It inverts where an opponent's weakness is temporary and finishing them is truly cheaper than meeting them re-armed --- the law's own reversal: destroy a man entirely or leave him entirely alone, and know which the situation priced.
\end{conditions}
\begin{application}
  \item Write the target before the campaign, with a date. No written mark, no discipline.
  \item Define in advance what will count as enough --- the terms, the territory, the concession.
  \item When the yes arrives, end the meeting. Nothing said after the win is for the deal.
  \item In victory, hand the defeated a role or a clean exit --- never a slow bleed. Spectators are watching the arithmetic.
  \item Audit winning streaks: three in a row is when the re-aiming starts. Re-read your own written mark.
\end{application}
\begin{feedback}
  \item Landing: the defeated can describe the deal without bitterness. The line was clean.
  \item Landing: the next negotiation opens easily --- spectators concluded you keep your word and your lines.
  \item Failing: you are renegotiating your own agreed terms. The reveller is driving.
  \item Failing: new fronts are opening before the old one closed. Momentum has the wheel.
\end{feedback}
\begin{failure}
The discipline calcifies into timidity --- a person who always stops at the stated mark becomes predictable to anyone who can read the mark, and some victories are only safe if completed. The law's reversal concedes exactly this: the skill is pricing completion against the new enemies, not kneeling in front of either choice. And the line requires a planner; a practitioner who cannot write targets cannot have a goal line, only moods that happen to halt.
\end{failure}
\begin{countermeasures}
  \item When a counterpart keeps adding asks after the yes, stand and go --- you are watching victory intoxication from inside its blast radius.
  \item In your own house, watch the executive who wins publicly and litigates privately. The line was crossed inward, where it counts.
  \item Price your own streak: what did you agree to in the week after your last win that you would have refused the week before? (T-24-02's freeze applies to wins, not only gifts.)
  \item When someone urges you past your stated goal for your own good, ask whose goal it has become.
\end{countermeasures}

\sources{B3}
\seesources
\seealso{T-28-01, T-17-01, T-21-01}
